\section{Introduction}

Independently fine-tuned models built on one pre-trained checkpoint have task vectors whose
input-side singular subspaces overlap far above an isotropic random baseline, and a growing body
of merging work measures that overlap and builds on it.
\citet{gargiulo2025task} report the cross-task inner products of right singular vectors and
define an interference score from them.
\citet{demystify2026} define a Right Subspace Overlap statistic over $190$ pairs of task vectors
and use it as one input to a predictor of merging success.
\citet{qiu2026nufilt} construct a null-space projector out of the cumulative update's right
singular vectors, and justify it with a theorem that task vectors approximately align with
representation subspaces.
In each case the overlap is compared, explicitly or implicitly, against what an unstructured
subspace would give.

That comparison is correct for unconditional random subspaces, but it does not answer the
conditional attribution question studied here. For a linear map $y=Wh$, the raw minibatch
gradient has the outer-product form
\begin{equation}
 \nabla_W L_t=\sum_{b\in\mathcal B_t}g_{t,b}h_{t,b}^\top.
 \label{eq:outer-product}
\end{equation}
Under scalar-step gradient descent without weight decay, the final update's row space therefore
lies in the span of the activations along that specialist's own training trajectory. This exact statement neither
identifies a common frozen-base span nor applies unchanged after coordinatewise optimiser
preconditioning: activations drift and optimisers transform gradients. It motivates a weaker,
measurable question. Using the frozen base's activation second moment as a common coordinate
system, how much observed overlap is predicted after holding fixed each update's own occupancy
of that geometry? The isotropic value $k/d_{\mathrm{in}}$ conditions on none of it.

Prior work motivates both sides of this question. \citet{cheng2025wudi} relate linear-layer task
vectors to their inputs, and \citet{qiu2026nufilt} give an approximate subspace-alignment result;
the latter note that the effect of shared initialisation remains insufficiently explored. These
results do not establish frozen-base containment for the checkpoints studied here. Cross-task
right-subspace overlap itself is already measured at scale by the work above. What is missing is
a stated conditional baseline, an accounting of how much the answer moves under it, and evidence
that the baseline randomises what it claims to.

We provide all three. The activation-conditioned null independently acts on each observed update
with rotations that preserve its complete singular spectrum and blockwise read-mass profile in
the frozen-base activation eigenbasis. Its expectation has a closed form: it is the overlap
predicted by the two observed occupancy profiles. The remainder is residual directional
agreement conditional on them. This is an operational conditional decomposition; the retained
profile is measured after fine-tuning and may itself contain task-learned structure. A generative
alternative sampled from $C_H$ is not a valid conditional comparator because it must invent the
update--activation relationship and produces a conspicuous level failure in our
common-reference stress test.

Applied to the CLIP benchmarks used by this literature, the null reproduces $25.3\%$ and $22.4\%$
of above-isotropic top-$8$ overlap in the two complete ViT-B measurements. Conditional on the
fixed calibration and blocking specification, the remaining excess stays positive throughout the
crossed task--layer finite-pool stability interval. Under rank-$16$
LoRA, $\mathrm{row}(BA)\subseteq\mathrm{row}(A)$ makes the observed $A$ factor a necessary
regime-specific condition. Its top-$8$ right singular subspaces overlap at $0.88$, and at $1.00$
when linearisation keeps $A$ bit-identical, while the learned $B$ factor's top-$8$ left singular
subspaces overlap at $0.05$ against $0.01$ isotropic
chance. An $A$- and factor-spectrum-conditioned null reproduces $80.3\%$ of the product update's
above-isotropic overlap;
although its attainable plant delivers only $37.4$--$57.0\%$ of the nominal increment, the null
recovers $100.7$--$101.1\%$ of what is delivered. This one stress test does not identify which
mechanism produced the conditioned structure.

The methodological result is why these estimates, rather than a larger number from our first
implementation, are reported. A zero-signal level check accepts the identity map, so it cannot
detect under-randomisation. We add a positive control that plants shared directions and measures
the signal actually delivered and recovered. It exposed two under-randomisation defects in the
activation null and a nuisance-spectrum mismatch in the LoRA null. Correcting the activation
null raised delivered-signal recovery from $45.6$--$51.7\%$ to $99.3$--$99.6\%$ and moved the full-
fine-tuning ViT-B/16 headline from roughly $67\%$ to $25.3\%$. Surrogates independently drawn from that conditional
null also form a relative stress test: the isotropic and participation-ratio-whitened baselines
report $+0.0171$ and $+0.0522$ apparent excess. This comparison is conditional on our stated
reference model and is not an independent intervention.

We make three contributions: a group-conditional decomposition of read-subspace overlap; a
positive-control protocol that checks power as well as level and catches consequential defects in
our own nulls; and a regime-specific audit of full fine-tuning and LoRA that separates effect-size
attribution from pair ranking and downstream merge performance.
